\cotefile{156-8}
\scanpage{18}
\mspage{9}
% archive sheet number: 18
% lines 1--2
\[
\begin{aligned}
\text{b/}\quad &\unsure{\operatorname{envelopp}|\operatorname{envelop}}(F)\unsure{\ll|\subset\subset}\operatorname{envelopp}G\Longrightarrow F\ll G\\
\text{c/}\quad &\operatorname{envelopp}(F)\unsure{\gtrless|\bowtie}\operatorname{envelopp}(G)\Longrightarrow F\gtrless G.
\end{aligned}
\]
% lines 3--7
C'est quand ces \deleted{\unsure{conditions}} implications sont satisfaites \inserted{(en plus de la condition a) ci-dessus)} (et il suffit de le poser pour $F=X$, $\deleted{\unsure{F}}G=Y$, $X,Y\unsure{\in M}$ \inserted{(spécimens)}) qu'on dit que le \unsure{magasin} est ``ensembliste''.

\deleted{\unsure{Il est donc donné par un ensemble de}}

\deleted{Une telle donnée \unsure{équivaut|équival}}
% lines 8--14
\uline{Proposition}\inserted{-- \uline{Scholie}}. La donnée d'un \unsure{magasin} ensembliste, \unsure{à isom. près} revient à la donnée d'un ens. $L$ (l'ens des ``lieux'' ou ``points''), et d'une partie
\[
M\subset\operatorname{Figel}(L),
\]
satisfaisant les conditions
% lines 15--20
\uline{\unsure{Mag ens|Magas}} 1\quad $\forall\deleted{\unsure{\mathscr{F}}}F\in M$, et tout $\deleted{\unsure{x}}X\in\deleted{\unsure{\mathscr{F}}}F$, l'ens.
\[
\deleted{\unsure{\mathscr{F}}}F_{\unsure{\leq X|\subset X|X}}=(Y\in F\mid Y\subset X)\subset F
\]
est aussi dans $M$ ($M$ stable par \unsure{passage aux sous-multistrates}).

\uline{\unsure{Mag ens|Magas}} 2\quad $\forall x\in L$, $\{\{x\}\}\in M$.
% slanted
\margin{!!! Donc on a une \unsure{espèce de structure} $M\in\unsure{\mathscr{P}|\mathcal{P}}\mathscr{P}(L)$, avec \unsure{axiomes} \uline{\unsure{Mag ens|Magas}} 0 : $\forall X\in M$ est une figure ensembliste élémentaire, -- plus les deux \unsure{autres}, à peine plus \unsure{compliqué} que la définition d'un \unsure{espace} topologique !}
% lines 21--28
\unsure{Ça} c'est donc pas mal dégonflé !!!

En fait, à ce stade de l'axiomatique \unsure{``naïviste : là''|``naïviste'' là}, le point de vue ``atlas'' des figures nous avait plutôt compliqué la vie. Ici il vaut décidément mieux partir du \unsure{magasin}, et y voir aussi clair que possible -- quitte à faire ensuite une \unsure{place} \unsure{convenable|amplement} \unsure{aux|sur les} atlas plus
\scanpage{19}
\mspage{10}
% archive sheet number: 19
% lines 1--5
généraux que l'atlas canonique associé à un magasin -- et de revoir les prop. essentielles prop 1 à prop 5 (plus d'autres à venir peut-être) dans cet angle, \unsure{logiquement} \unsure{effleuré|effleurant|effarant} à peine.

% line 6
\uline{27 juin} (ou plutôt, 28 juin ! \unsure{l'on ne sait plus|on ne sait plus|(on ne sait plus}\ldots)

% lines 7--11
\uline{Définition}. Soit $\unsure{\mathscr{F}|\mathcal{F}},\leq$ un ensemble ordonné. On appelle \uline{antifiltre} de $\mathscr{F}$ une partie $\mathscr{F}'$ de $\mathscr{F}$ satisfaisant les conditions suivantes

a) $\mathscr{F}'$ est fermé pour $\leq$, i.e. $F\in\mathscr{F}'$, $G\leq F\Rightarrow G\in\mathscr{F}'$.

b) Si $F,G\in\mathscr{F}'$, et si $F\vee G$ existe dans $\mathscr{F}$, alors $F\vee G\in\mathscr{F}'$.

% lines 12--13, cancelled
\deleted{On dit que l'antifiltre est \uline{\unsure{couvrant}} si}
\[
\deleted{\unsure{\bigcup_{F\in\mathscr{F}}F=\Omega|\bigcup_{F\in\mathscr{F}}F=\mathscr{F}}}
\]

% lines 14--20
\uline{Définition}. Soit $M$ un magasin. Un \uline{atlas \unsure{attaché|associé} à $M$}, \inserted{ou $M$-atlas,} est par définition un antifiltre $\mathscr{F}$ (dans l'``ens'' $\mathscr{F}_M$ des \inserted{$M$-}figures, \unsure{inconditionnelles}, \deleted{tel que la figure vide et} contenant les \deleted{figure vide} figures élémentaires $\deleted{\unsure{X}}\widetilde{X}$ $(X\in M)$ et la figure vide \deleted{Cette dernière condition résulte des autres si $M\neq\varnothing$}).

% margin note written sideways
\margin{\uline{\unsure{Mag atl}} \inserted{sur $\mathscr{F}\subset\mathscr{F}_M$}\par
a) $F\in\mathscr{F}$, $G\leq F\Rightarrow G\in\mathscr{F}$.\par
b) $F,G\in\mathscr{F}$, $F\unsure{\gtrless|\bowtie}G\Rightarrow F\cup G\in\mathscr{F}$.\par
c) $\forall X\in M$, $\widetilde{X}\in\mathscr{F}$, et\par
d) $\mathscr{F}\neq\varnothing$ (i.e. $\varnothing\in\mathscr{F}$).}
% margin note written sideways
\margin{\unsure{revient à dire que} $\mathscr{F}$ (en tant que $\subset\unsure{\mathscr{P}|\mathcal{P}}(M)$) \uline{couvre} $M$ $\Longleftrightarrow$ c)}

% lines 21--26
\uline{NB} Cette dernière condition résulte des autres si $M\neq\varnothing$. On voit que les conditions suivantes sont équivalentes
\[
\begin{aligned}
\text{(i)}\quad&M=\varnothing\\
\text{(ii)}\quad&\mathscr{F}=\{\varnothing\}\\
\text{(iii)}\quad&\operatorname{Card}\mathscr{F}=1
\end{aligned}
\]
% lines 27--30
On dit \deleted{que} alors que \uline{le magasin} \deleted{$M$ est} \inserted{$M$ est} \uline{vide}, ou que \uline{l'atlas} \deleted{$\mathscr{F}$} \inserted{$\mathscr{F}$ est} \uline{vide} (un abus de langage pour $\mathscr{F}$, car $\mathscr{F}$ en tant qu'ens n'est pas vide \ldots.
\scanpage{20}
\mspage{11}
% archive sheet number: 20
% lines 1--4
Les \deleted{structures}\inserted{relations} $\leq$, $\ll$, $\unsure{\gtrless|\bowtie}$, $\parallel$ sur $\unsure{\mathscr{F}|\mathcal{F}}_{M^{\unsure{d|a|\circ}}}$ induisent des \deleted{$\unsure{\Rightarrow}$}\deleted{structures}\inserted{relations} de même nom sur $\mathscr{F}$.
Les prop $\unsure{1_{\circ}\mathrm{F}|1^{\circ}\mathrm{F}|1_0\mathrm{F}|1_o\mathrm{F}}$, $2_{\circ}\mathrm{F}$ (p. 2 bis), prop. 2, 3, $\deleted{\unsure{3}}5$ \unsure{pour|passent} ``ipso facto'' \unsure{se conservent} dans un atlas général.

% lines 5--8
Il n'y a que la prop. 4 (recollement des raffinements) qui, dans le cas d'un ens d'indices ``\uline{infini}'' $I$, demande un axiome de stabilité supplémentaire pour l'atlas,

% lines 9--14
Savoir\par
\uline{At Fil}\quad Soit $F\in\mathscr{F}$, et soit $(\deleted{\unsure{F|E}}G_\alpha)$ une famille filtrante croissante \inserted{pour $\leq$} de raffinements de $F$. Supposons que $\forall X\in\widetilde{F}$, la famille des $G_\alpha\mid X$ soit stationnaire. Alors \deleted{$\unsure{\operatorname{Sup}G_\alpha|\bigcup G_\alpha}$} $\operatorname{Sup}^{\leq}G_\alpha$ existe dans $\mathscr{F}$ \deleted{(et est un raffinement} \deleted{\unsure{(i.e.}}.

% lines 15--18
Du point de vue $\mathscr{F}\subset\mathscr{F}_M$, cela signifie que $\bigcup G_\alpha$ (qui est le Sup dans $\mathscr{F}_M$) est en fait dans $\mathscr{F}$ : propriété de stabilité de $\mathscr{F}$ en \deleted{\unsure{tant}} que partie fermée de $\mathscr{F}_M,\leq$.

% lines 19--22
On peut considérer la structure d'atlas comme donnée indépendamment de celle d'un magasin donné au préalable.

Atlas : $\mathscr{F}$ muni de trois relations $\leq$, $\ll$, $\unsure{\mid\circ\mid|\parallel}$

% lines 23--27
\uline{At 1}\quad Les Sup majorés existent, et $\mathscr{F}\neq\varnothing$ ($\Rightarrow\mathscr{F}$ a un plus petit élément $\varnothing_{\mathscr{F}}$).

\deleted{\unsure{Éléments}} Figures élémentaires \unsure{$\cup$|$\vee$}-éléments strictement irréductibles :
\[
X=\operatorname{Sup}F_i\Rightarrow\exists i,\ \deleted{\unsure{K_i}}F_i=X.
\]
$\deleted{\unsure{E}}M$ ens des fig. élém., $\widetilde{F}=\{X\in M\mid X\leq F\}$.

% slanted
\margin{At 2, 3, 4 : \uline{passage aux strates}}
% lines 28--30
\[
\text{\uline{At 2}}\quad F\leq G\Longleftrightarrow\widetilde{\deleted{\unsure{f}}F}\subset\widetilde{G}\qquad\deleted{\unsure{(i.e.)}}\text{i.e. }X\in M,\ X\leq F\Rightarrow X\leq G.
\]
``toute strate de $F$ est \unsure{strate} de $G$''.

\unsure{Equivalent|Équivalent} : Toute figure est le Sup de ses strates
\scanpage{21}
\mspage{12}
% archive sheet number: 21
% slanted
\margin{\uline{NB} At 1--At 3 ne concernent que $\leq$. At 4 ne fait que préciser le \unsure{passage} de $\ll$ aux strates}
% lines 1--4
\[
\text{Déf}\quad F\unsure{\gtrless|\bowtie}G\Longleftrightarrow\{F,G\}\text{ majoré }(\Longleftrightarrow F\unsure{\cup|\vee}G\text{ existe})
\]
\[
\text{\uline{At 3}}\quad F\gtrless G\Longleftrightarrow\forall X\in\widetilde{\deleted{\unsure{X|F}}F},\ Y\in\widetilde{G},\ X\gtrless Y
\]
(NB $\Rightarrow$ est déjà connu)
\[
\text{\uline{At 4}}\quad F\ll G\Longleftrightarrow\forall X\in\widetilde{F},\ \exists Y\in\widetilde{G},\ X\ll Y
\]
% lines 5--7
\[
\text{\uline{At 5}}\quad\text{a)}\quad F\mid\circ\mid G\Longrightarrow F,G\in M.
\]
\margin{i.e. la relation $\mid\circ\mid$ est en fait une relation dans $M$}
% slanted
\margin{\unsure{relie} $\mid\circ\mid$ et $\gtrless$}
\[
\begin{aligned}
\text{b)}\quad X\deleted{\unsure{\mid\circ\mid}}\gtrless Y\Longleftrightarrow{}&\forall\unsure{X|X'}\in\widetilde{X}\setminus\widetilde{X}\cap\widetilde{Y},\ Y''\in\widetilde{Y}\setminus\widetilde{X}\cap\widetilde{Y},\;\unsure{\text{--}}\\
&\unsure{\text{--}}\quad X'\mid\circ\mid Y''
\end{aligned}
\]
\[
+\operatorname{Mag}1,\operatorname{Mag}2,\operatorname{Mag}3,\operatorname{Mag}4
\]

% lines 8--14
\deleted{\illegible{1}}
\uline{Scholie} Les deux pts de vue
\[
\begin{aligned}
\text{(i)}\quad &(M,\leq,\ll,\mid\circ\mid,\unsure{\mathscr{F}|\mathcal{F}}\subset\unsure{\mathscr{F}_M|\mathbb{F}_M|F_M|\mathscr{F}_{M^{\zeta}}|\mathscr{F}_{M^{\varphi}}|\mathscr{F}_{M^f}})\quad\text{avec}\quad
\underbrace{\operatorname{Mag}1-\operatorname{Mag}4}_{\substack{\text{concernant}\\\leq,\ll,\mid\circ\mid}}+\\
&\operatorname{Mag\,at}\\
\text{(ii)}\quad &(\mathscr{F},\leq,\ll,\mid\circ\mid)\quad\text{avec axiomes}\quad\operatorname{At}1-\operatorname{At}5\\
&+\operatorname{Mag}1-\operatorname{Mag}4
\end{aligned}
\]
sont équivalents, i.e. donnent des catégories-groupoïdes \unsure{équivalents|équivalentes}.

% lines 15--22
Le point de vue ``magasin'' est plus pratique, visiblement, pour l'énoncé et le \unsure{dévissage} des axiomes de base.
Le point de vue ``atlas'' me semble par contre plus commode pour la pratique du calcul et des raisonnements géom, une fois l'axiomatique acquise.
\scanpage{22}
\mspage{13}
% archive sheet number: 22
% lines 1--5
Point de vue ``symétrique''. Structure sur deux ensembles de base $M$ (multistrates), $\unsure{F|\mathscr{F}}$ (figures)

Structure\par
a) Relation d'incidence $X\triangleleft F$ entre $M,F$ (suffira à définir $\leq$ sur $M$ et sur $F$)

b) Relations $\ll$, $\mid\circ\mid$ sur $M$

% lines 6--7
Pour tout $F\in\unsure{\mathscr{F}|\mathcal{F}}$, on pose
\[
\widetilde{F}=\{X\in M\mid X\triangleleft F\}\qquad
\text{On pose }F\leq G\overset{\text{déf}}{\Longleftrightarrow}\widetilde{F}\subset\widetilde{G}
\]

% lines 8--14: entire draft block crossed by long diagonals
\margin{Ainsi,}
\deleted{\uline{At M1} $\widetilde{F}=\widetilde{G}\unsure{\Rightarrow}\Rightarrow F=G$, i.e. $\mathscr{F}\longrightarrow\unsure{\mathscr{P}|\mathcal{P}}(M)$, $F\longmapsto\widetilde{F}$ est \unsure{injection|injective}, ou encore le préordre $\leq$ sur $\mathscr{F}$ est un \uuline{ordre}.}

\deleted{\uline{At M2} $\unsure{F}\ \widetilde{\mathscr{F}}$ \unsure{couvre} est une \unsure{figure élémentaire dans}\inserted{famille de parties de $M$ quasi-admissible} $\unsure{\mathscr{F}}$ qui couvre $M$ et sépare les points de $M$, i.e.}

\deleted{a) $\forall X\in M$, $\unsure{\mathscr{F}}\ \mathscr{F}^{X}=\{F\in\mathscr{F}\mid X\triangleleft F\}$ a un plus petit élément (on l'appelle la \uline{figure élémentaire} ou multistrate déterminée par $X$)}

\deleted{\unsure{b)}}

% lines 15--20: revised axioms
\uline{At M1} La relation $X\triangleleft F$ est séparante en $M$ et en $\mathscr{F}$, i.e.

a) L'application $F\longmapsto\widetilde{F}$, $\mathscr{F}\longrightarrow\mathscr{P}(M)$ est injective [ou encore, $\leq$ est une relation d'\uline{ordre}]

b) $\deleted{\unsure{\forall}}$ Si $X,Y\in M$ tels que $X\neq Y$, alors $\exists F\in\mathscr{F}$ avec $X\in\widetilde{F}$, $Y\notin\widetilde{F}$ ou l'inverse (i.e. $\widetilde{\mathscr{F}}\subset\mathscr{P}(M)$ sépare les points de $M$)

% lines 21--24
\uline{At M2} La famille de parties $\widetilde{\mathscr{F}}$ dans $M$ est quasi-admissible, i.e. $\forall X\in M$, $\deleted{\unsure{\mathscr{F}}}$ posant $\mathscr{F}^{X}=\{F\in\mathscr{F}\mid X\triangleleft F\}$, $\mathscr{F}^{X}$ a un plus petit élément (appelé la \uline{figure élémentaire} associée à la multistrate $X$)

% lines 25--29
Ainsi \deleted{$X\triangleleft F$ équivaut à}
\[
X\triangleleft F\Longleftrightarrow F_X\leq F\qquad(\Longleftrightarrow X\in\widetilde{F})
\]
de sorte que la condition At M1 b), compte tenu de At M1 a), équivaut à
\[
M\longrightarrow\mathscr{F}\quad X\longrightarrow F_X\qquad\text{est injectif}.
\]

% lines 30--32
\uline{At M3} a) Soit $\deleted{\unsure{\mathscr{F}}}\ \unsure{\underline{\Phi}_i|\underline{\varphi}_i},\underline{\Phi}\in\widetilde{\mathscr{F}}$, avec $\underline{\Phi}_i\subset\underline{\Phi}$. Alors $\bigcup\underline{\Phi}_i\in\widetilde{\mathscr{F}}$ (stabilité de $\widetilde{\mathscr{F}}$ par $\unsure{\bigcup_{\alpha}|\bigcup}$ \unsure{majorées})

\margin{résulte du reste si $M\neq\varnothing$}
b) $\unsure{\widetilde{\mathscr{F}}|\mathscr{F}}\neq\varnothing$ ($\deleted{\Longleftrightarrow}\varnothing\in\widetilde{\mathscr{F}}\Longleftrightarrow\widetilde{\mathscr{F}}\neq\varnothing\Longleftrightarrow\mathscr{F}$ a un plus petit élément)
\scanpage{23}
\mspage{14}
% archive sheet number: 23
% lines 1--3
Ces conditions At M1--At M3 \inserted{plus At M4 plus bas} ne concernent que \deleted{$\unsure{\Leftarrow|\leftarrow|\triangleleft}$} la relation $X\triangleleft F$. Allons un peu plus loin. Si $X,Y\in M$, on a tautologiquement
% lines 4--6
\[
\begin{aligned}
F_X\leq F_Y&\Longleftrightarrow\forall\deleted{\unsure{\underline{\Phi}\in\widetilde{\mathscr{F}}}}F\in\unsure{\mathscr{F}|\mathcal{F}},\ \deleted{\unsure{X}}Y\triangleleft F\Rightarrow X\triangleleft F\\
&\Longleftrightarrow\forall\unsure{\underline{\Phi}|\underline{\varphi}}\in\widetilde{\mathscr{F}},\ Y\in\underline{\Phi}\Rightarrow X\in\underline{\Phi}\\
&\Longleftrightarrow X\triangleleft F_Y\Longleftrightarrow X\in\widetilde{F}_Y.
\end{aligned}
\]
% lines 7--11
On écrit alors $X\leq Y$, c'est une relation de préordre, savoir la relation de préordre associée à la famille \unsure{préadmissible|quasi-admissible} $\widetilde{\mathscr{F}}$ dans $M$. Par At M1 b) c'est une relation d'\uline{ordre}. Ceci dit, on a une
% slanted
\margin{On écrit $\widetilde{X}=\widetilde{F}_X=\{Y\mid Y\triangleleft\deleted{\unsure{f}}F_X\}=\deleted{\unsure{\mathscr{F}}}M_{\leq X}$}

% lines 12--14, cancelled proposition
\deleted{\uline{Prop} Si $F\in\mathscr{F}$, $\inserted{\underline{\Phi}=}\widetilde{F}$ est une partie fermée de $M,\leq$.}
\deleted{Si $\unsure{C|E}\ \Psi$ est une partie fermée de $\underline{\Phi}$, alors $\Psi\in\widetilde{\mathscr{F}}$}

% lines 15--16
\uline{Prop} Soit $\underline{\Phi}\in\unsure{\mathscr{F}|\widetilde{\mathscr{F}}}$, alors toute partie fermée \inserted{$\Psi$} de $\underline{\Phi}$ est dans $\widetilde{\mathscr{F}}$.

% lines 17--23
Ceci \deleted{\unsure{renforce At}} est formellement plus fort que At M3 a), et lui est donc équivalent (compte tenu du reste). Si j'ai laissé At M3 a) sous forme apparemment plus faible, c'est pour lui donner un sens en termes de $(\mathscr{F},\leq)$, comme axiome d'existence de Sup majorés.

% lines 24--28
\deleted{\uline{At M4} Soit $F,G\in\mathscr{F}$. Pour que $\unsure{F}$}

Soit $F,G\in\mathscr{F}$, on pose
\[
\begin{aligned}
F\unsure{\gtrless|\bowtie}G&\overset{\text{déf}}{\Longleftrightarrow}\{F,G\}\text{ majoré dans }\mathscr{F}\Longleftrightarrow F\cup G\text{ existe}\\
&\Longleftrightarrow\widetilde{F}\cup\widetilde{G}\in\widetilde{\mathscr{F}}
\end{aligned}
\]
\[
X\gtrless Y\overset{\text{déf}}{\Longleftrightarrow}F_X\gtrless F_Y\quad\text{i.e.}\quad\widetilde{X}\cup\widetilde{Y}\in\widetilde{\mathscr{F}}\quad[\widetilde{X}=M_{\leq X}]
\]
\[
\text{\uline{At M4}}\quad F\gtrless G\Longleftrightarrow\forall X\in\widetilde{F},\ Y\in\widetilde{G},\text{ on a }X\gtrless Y
\]
% lines 29--30
\deleted{\unsure{At}} Pour les relations \unsure{données} $\ll$, $\unsure{\circ\mid|\circ|\mid\circ\mid}$ sur $M$, on pose
\[
\text{\uline{At M5}}\quad X\gtrless\deleted{\unsure{F|G}}Y\Longleftrightarrow\forall X'\in\widetilde{X}\setminus\widetilde{X}\cap\widetilde{Y},\ Y'\in\widetilde{X}\cap\widetilde{Y},\text{ on a }X'\mid\circ\mid Y'.
\]
\scanpage{24}
\mspage{15}
% archive sheet number: 24
% lines 1--2
\uuline{Plus} les axiomes \inserted{des magasins} reliant \uline{$\leq$, $\ll$, $\circ$ sur $M$}, savoir Mag 1--Mag 4.

% lines 3--8
Ainsi, qu'on parte de $(\unsure{\mathscr{F}|\mathcal{F}},\leq,\ll,\mid\circ\mid)$, ou de $(\mathscr{F},M,\triangleleft,\ll,\mid\circ\mid)$, il nous faut, en plus de Mag 1--Mag 4, cinq autres axiomes.
Si on part de $(M,\leq,\ll,\mid\circ\mid,\mathscr{F})$, il nous faut un seul axiome (à tiroirs\ldots) supplémentaire, concernant le choix de $\mathscr{F}$.

% lines 9--13
Tout ceci concerne les atlas que je devrais appeler ``\unsure{spaciaux|spatiaux|spéciaux}''. Il faut maintenant ouvrir une parenthèse sur les \unsure{``atlas ``polyédraux''|``atlas'' ``polyédraux''}, qui représentent une \inserted{espèce de} structure très légèrement différente.

% lines 14--17
\uline{Définition} Le magasin $M$ est dit \inserted{pré}\uline{polyédral} si \deleted{$\forall X\in M$, $\widetilde{X}$ est un ens. ordonné polyédral, i.e.}\inserted{$(M,\leq)$ est un ens. ordonné} polyédral, i.e. si deux éléments $X,Y$ de $M$ qui sont majorés et minorés, ont un Inf \inserted{dans $M$}.

% lines 18--21: cancelled block, including its interlinear revision
\deleted{\inserted{Il est dit \uline{polyédral} si de plus, cette propriété est}}
\deleted{(\uline{NB} Si les $M_{a\leq\cdot\leq b}$ sont finis, \unsure{cette exigence|c'est l'exigence} est \unsure{entraînée|entendue} -- plus \unsure{gén.}, si $M$ est polyédral et si les $M_{a\leq\cdot\leq b}$ sont \uuline{inductifs}, alors \unsure{deux} \unsure{max. éléments}}

% lines 22--26
On introduit \deleted{\unsure{donc}} les relations plus fortes que $\ll$ (pour les magasins, pas néc. prépolyédraux)
\[
\begin{aligned}
X\underset{\mathrm{pol}}{\overset{\unsure{\mathrm{pr}|\mathrm{pré}}}{\ll}}Y
\Longleftrightarrow{}&X\ll Y,\ \text{\uuline{et}}\ \forall Y'\in\widetilde{Y},\quad
\widetilde{X}_{Y'}\overset{\text{déf}}{=}\widetilde{X}\cap\operatorname{Oubl}(\unsure{Y|Y'}
% right edge clips the Oubl argument; no closing parenthesis is supplied.
\end{aligned}
\]
est vide \uline{ou} a un plus grand élément, (i.e. $X_{Y'}$ est une figure élémentaire)

% bottom boxed margin note
\margin{En fait $X_{Y'}=X\unsure{\bullet|\cdot}Y'=\operatorname{Inf}^{\ll}(X,Y')$ dans $\mathscr{F}_M$, dans le formalisme des figures\ldots}
\scanpage{25}
\mspage{16}
% archive sheet number: 25
% lines 1--3
La condition que $M$ est \inserted{pré}polyédral \deleted{\unsure{est ici}} équivaut à dire que
\[
X\leq Y\Longrightarrow X\ll_{\mathrm{pol}}Y
\]
% lines 4--6
ce qui est nécessaire pour que
\[
(M,\leq,\ll_{\mathrm{pol}},\mid\circ\mid)
\]
soit encore un magasin. Est-ce suffisant ?

% lines 7--9
1°) Mag 1\quad Il reste à vérifier que $\ll_{\mathrm{pol}}$ est une relation d'ordre, et tout revient à la transitivité
\[
X\ll_{\mathrm{pol}}Y\ll_{\mathrm{pol}}Z\Longrightarrow X\ll_{\mathrm{pol}}Z.
\]
% lines 10--16
Soit $Z'\leq Z$, alors par $Y\ll_{\mathrm{pol}}Z$ \deleted{\unsure{donc}} la sous-figure induite $\unsure{Y^{4}_{Z'}|Y^{*}_{Z'}|Y^{+}_{Z'}|Y_{Z'}}$ est vide, ou de la forme $Y'$. Dans le premier cas, la figure induite par $\unsure{Y_{Z'}|Z'}$ sur $X$ est vide, OK, dans le deuxième elle est vide ou de la forme $X'$, par $\ll_{\mathrm{pol}}$, OK encore. (On utilise transitivité des images inverses)

% lines 17--22
2°) \uline{Mag 2}\quad Soit $X\ll_{\mathrm{pol}}Y$, prouvons que l'ens
\[
\widetilde{Y}^{X,\mathrm{pol}}=\{Z\in\widetilde{Y}\mid X\ll_{\mathrm{pol}}Z\}
\]
a un plus petit élément, pour ceci, il suffit de noter que l'inclusion évidente $\subset$ est une égalité
\[
\widetilde{Y}^{X,\mathrm{pol}}=\widetilde{Y}^{X},
\]
\deleted{est en fait une égalité},
% lines 23--26
i.e. que \deleted{$\unsure{X}$ \unsure{Si|$S'$}}
\[
X\ll Y'\leq Y\Longrightarrow X\ll_{\mathrm{pol}}Y',
\]
c'est évident par définition (vu que les $Y''\in\widetilde{Y}'$ sont aussi $\in\widetilde{Y}$).

% lines 27--28
\deleted{\unsure{3°)}} \uline{NB} On en conclut que
\[
X\overset{\circ}{\ll}_{\mathrm{pol}}Y\Longleftrightarrow X\ll_{\mathrm{pol}}Y\text{ et }X\overset{\circ}{\ll}Y
\]
